\textbf{Задача 20} Докажите, что язык \textsf{MATRIXMDEG} = $\{(A \in \{0, 1\}^{n \times n}, 1^k, i, j) | (A^{\lor k})_{ij} = 1 \}$ лежит в $\textbf{NC}$. Используя этот факт, докажите, что $\textbf{NL} \subset \textbf{NC}$.

\begin{proof} Достаточно доказать, что $\mathsf{PATH} \in \textbf{AC}^1 $, а логарифмическая сводимость не выводит за пределы этого класса. Начнём с того, почему $\mathsf{PATH} \in \textbf{AC}^1$. Пусть ориентированный граф с $N$ вершинами задан матрицей смежности $A$. Будем считать, что на диагонали стоят единицы, т.е. $A_{ij} = 1$, если из $i$ в $j$ есть путь длины не больше 1. Рассмотрим $\lor$-умножение булевых матриц, отличающееся от обычного умножения тем, что вместо сложения используется дизъюнкция: $(A \cdot B)_{ij} = \lor_{k=1}^{N}(A_{ik} \land B_{kj})$.
Тогда $(A^m)_{ij} = 1$, из $i$ в $j$ есть путь длины не больше $m$. (Возведение в степень, разумеется, понимается как итерированное $\lor$-умножение). Соответственно, $(A^N)_{ij} = 1$, если есть хоть какой-то путь из $i$ в $j$: если путь есть, то есть и путь не длиннее $N$. Таким образом, для решения задачи о наличии пути нужно возвести булеву матрицу в $N$-ую $\lor$-степень. Это нужно слелать быстрым возведением: посчитать $A^2$, затем $A^4$, $A^8$ и т.д., всего $\log N$ возведений в
квадрат. Каждый этап вычисляется схемой константной глубины (с элементами произвольной входящей степени) по определению. Таким образом, общая глубина схемы составит $O(\log N)$.

Теперь докажем, что логарифмическая сводимость не выводит из класса $\textbf{AC}^1$. На самом деле достаточно доказать это не для произвольной сходимости, а для конструкции, которая используется для сведения к $\mathsf{PATH}$. Иными словами, нужно установить, соответствует ли переход от одной конфигурации
к другой командам некоторой недетерминированной машины Тьюринга. Это делается уже известным из теоремы Кука–Левина методом: нужно, чтобы во всех малых окрестностях переход был допустимым. Такая формула будет конъюнкцией формул фиксированного размера, т.е. лежать в $\textbf{AC}^0$. Таким образом,
мы получили даже $\textbf{AC}^0$-сводимость, которая по доказанному ранее будет логарифмической.

Осталось заметить, что размер полученного графа будет полиномиальным
от размера исходного входа: $N = poly(n)$. Таким образом, глубина итоговой схемы составляет $O(\log N) = O(\log n)$, что и требовалось.

Таким образом, мы доказали $\textbf{NL} \subset \textbf{AC}^1 \subset \textbf{NC}$, и теперь можем доказать общее утверждение о том, что логарифмическая сводимость сохраняет класс $\textbf{AC}^1$. Действительно, при логарифмической сводимости каждый бит результата вычисляется на логарифмической памяти и, как следствие, $\textbf{AC}^1$-схемой. Значит,
и все биты будут вычислены $\textbf{AC}^1$-схемой.

\end{proof}